Assign the $N$ variables in different ways to the rows of the Young diagram (their placement among the cells of any one row

\SourcePage{290}{293}
is immaterial). This procedure gives a collection of functions that are carried into one another by an arbitrary permutation of the variables\footnote{The order of symmetrization and alternation could instead be reversed: first alternate over the variables in each column, and then symmetrize over those in the rows. This produces nothing new, however, because the functions obtained by the two procedures are linear combinations of one another.}. It must nevertheless be stressed that these functions are not all linearly independent. In general, the number of independent functions is smaller than the number of possible assignments of the variables to the rows of the diagram; we shall not examine this point further here\footnote{A set of independent functions carried into one another forms a basis for an irreducible representation of the permutation group. The number of these functions is the dimension of the representation. For particles of spin $1/2$, it is found in Problem~1 of this section.}.
